% Required in the paper preamble for the shaded report excerpts:
% \usepackage{graphicx}
% \usepackage{subcaption}
% \usepackage[most]{tcolorbox}

\subsection{Diagnosing a completed fit}

After parameter estimation has finished, examine the result with:

\begin{verbatim}
pfit diagnose sessions/<session-name> <run-id>
\end{verbatim}

The run identifier is the name of a directory under
\texttt{sessions/<session-name>/outputs/}. If it is omitted, the most recent
completed run is selected.

Diagnosis is based on the saved run snapshot: the YAML configuration, data
files, generated JAX script, generated source model, fitted parameters, and
runtime provenance copied when the fit was run. It does not depend on the
current working files in the session. A previous result can therefore still be
diagnosed after the working YAML, data, or generated code have been changed.

For each dataset, the command recreates the fitted trajectory and compares it
with the measurements. In the running example, separate plots are produced for
the pulse and washout experiments. These show the measured signal together with
the fitted \texttt{signal\_model}. Residual values are still written to CSV
files and summarized in the diagnosis report, but they are not drawn as a
separate subplot in \texttt{fit\_expN.png}.

The diagnosis reports residual summaries, including root-mean-square error,
mean absolute error, bias, largest absolute residual, and simple time-pattern
evidence. Missing measurements are excluded from residual statistics, but
nonfinite model predictions are treated as simulation failures rather than
missing data. The report also checks whether fitted parameters are close to
their bounds, whether optimizer logs suggest that improvement was still
occurring, and whether the run used a full fit or a gradient-only restart.

If post-fit sloppiness analysis was performed during \texttt{pfit run}, the
diagnosis includes its status and summary. The diagnosis command does not
normally repeat the sloppiness calculation.

The main files produced are:

\begin{verbatim}
fit_diagnosis.txt
scientific_diagnosis.json
fit_expN.png
residual_expN_channelM.csv
\end{verbatim}

Here, \texttt{N} identifies the experiment and \texttt{M} identifies the
measured output channel. The text report is the most convenient place to begin,
while the plots and residual CSV files allow the user to inspect individual
datasets and output channels.

For the running pulse/washout example, representative diagnostic plots are shown
in Figure~\ref{fig:pfit-drug-response-diagnosis}. These figures were generated
from synthetic data created from the same model, with small added measurement
noise, followed by the ordinary \texttt{pfit run} and \texttt{pfit diagnose}
workflow.

\begin{figure}[t]
  \centering
  \begin{subfigure}[t]{0.48\linewidth}
    \centering
    \includegraphics[width=\linewidth]{figures/fit_exp1.png}
    \caption{Pulse experiment.}
  \end{subfigure}
  \hfill
  \begin{subfigure}[t]{0.48\linewidth}
    \centering
    \includegraphics[width=\linewidth]{figures/fit_exp2.png}
    \caption{Washout experiment.}
  \end{subfigure}

  \vspace{0.6em}
  \begin{subfigure}[t]{0.42\linewidth}
    \centering
    \includegraphics[width=\linewidth]{figures/sloppiness_spectrum.png}
    \caption{Post-fit sloppiness spectrum.}
  \end{subfigure}
  \caption{Example diagnostic outputs for the pulse/washout running example.}
  \label{fig:pfit-drug-response-diagnosis}
\end{figure}

A representative excerpt from \texttt{fit\_diagnosis.txt} for this example is:

\begin{tcblisting}{
  listing only,
  breakable,
  colback=gray!6,
  colframe=gray!35,
  boxrule=0.4pt,
  arc=1mm,
  left=2mm,
  right=2mm,
  top=1mm,
  bottom=1mm
}
Fit diagnosis

Artifacts:
- final_design_point.csv: present
- result_solution_exp1.csv: present
- result_solution_exp2.csv: present
- de_fitting.log: present
- NODE_fitting.log: present
- fitting_error.txt: missing

Optimizer summary:
- DE: 120 iterations, first loss 8.0936E-02,
      best loss 5.7755E-03, final loss 5.7755E-03
- NODE: 21 iterations, first loss 5.7754E-03,
        best loss 5.7753E-03, final loss 5.7753E-03

Experiments: 2
Aggregation: equal_experiment_mean
- Experiment 1: pulse.csv; loss=0.005219267638936577;
  output=result_solution_exp1.csv
- Experiment 2: washout.csv; loss=0.006331386422963499;
  output=result_solution_exp2.csv

Scientific diagnosis:
- Status: ok
- Refinement stopped at its iteration budget.
- [BUDGET] Refinement stopped at its iteration budget.
  Next: Assess the final gradient and trajectory plots before
  deciding whether a gradient-only restart is worthwhile.
- Figure: fit_exp1.png
- Figure: fit_exp2.png
- Detailed evidence: scientific_diagnosis.json
\end{tcblisting}

When sloppiness analysis is enabled, a representative
\texttt{sloppiness\_report.txt} excerpt is:

\begin{tcblisting}{
  listing only,
  breakable,
  colback=gray!6,
  colframe=gray!35,
  boxrule=0.4pt,
  arc=1mm,
  left=2mm,
  right=2mm,
  top=1mm,
  bottom=1mm
}
Post-fit sloppiness
Status: ok
Method: finite-difference
Coordinates: log10(k_elim), log10(k_on), log10(k_off), baseline, gain
Verdict: no weak curvature directions at the specified relative floor
Weak modes: 0
Negative modes: 0
Normalized gradient infinity norm: 5.63383e-06
Eigenvalues (largest to smallest):
  mode 1: 1.68783850e+01
  mode 2: 1.26587665e+00
  mode 3: 8.25273064e-02
  mode 4: 1.03427087e-03
  mode 5: 2.01732936e-04
Positive spectrum span: 4.92255 decades
STIFFEST direction: k_off(-0.64), k_elim(-0.63),
                   k_on(+0.43), baseline(+0.02)
SLOPPIEST direction: gain(+1.00), baseline(-0.08),
                    k_on(-0.02), k_elim(+0.01)
Warning: Local curvature only: not proof of structural/global
identifiability or convergence.
\end{tcblisting}

By default, the framework may ask the configured local language model to
summarize the computed evidence and suggest next steps. This interpretation is
advisory: the model receives the numerical diagnosis, but it does not modify the
session or the run. To produce only the deterministic report without contacting
the local model, use:

\begin{verbatim}
pfit diagnose sessions/<session-name> <run-id> --deterministic-only
\end{verbatim}

When a closer examination of local convergence is needed, the optional
\texttt{--probe-gradients} flag compares automatic derivatives with bounded
finite-difference estimates for a small selection of parameter axes:

\begin{verbatim}
pfit diagnose sessions/<session-name> <run-id> --probe-gradients
\end{verbatim}

This check can be computationally expensive and is not required for routine
diagnosis.

The diagnosis should be treated as evidence for scientific review, not as a
pass--fail judgment. A low loss does not prove that the fitted parameters are
identifiable, that the global optimum has been found, or that the model is
scientifically appropriate. The user should inspect the trajectories, residuals,
parameter locations, optimizer history, and sloppiness summary before deciding
whether the fit answers the intended scientific question.
